\section{Semantics of \egglog}

In this section we describe the semantics of core \egglog.
Core \egglog differs from the full \egglog language in several aspects.
For example,
 \egglog allows multiple actions in a rule 
 while core \egglog allows only one atom in the head,
 and core \egglog does not have the \lstinline{union} operation.
These \egglog features can be desugared into the core language.
However, there are also some assumptions we made about the core \egglog.
For example, we assume the \lstinline{:merge} expression over ids
 are \lstinline{union} and the \lstinline{:merge} expressions 
 over interpreted constants is the join operator of a given lattice,
 while \egglog allows \lstinline{:merge} to be any valid \egglog expression.
In other words,
 the core \egglog captures a well-behaving subset of the full \egglog language.

\subsection{Syntax}

Given the set of (interpreted) constants $C$,
 the syntax of the core \egglog language is shown in \autoref{fig:syntax}.
An \egglog program is defined as a list of rules, 
 and each rule consists of an atom in the head
 and a list of atoms in the body.
An atom has the form $f(p_1,\ldots,p_k)\mapsto o$ 
 and intuitively means function $f$ has value $o$ on $p_1,\ldots, p_k$.
A pattern $p$ is a nested expression 
 constructed using function symbols, variables, and constants.
We additionally define a ground term (or term) $t$ to be a pattern with no variables,
 and a ground atom to be an atom where all the patterns are ground terms.

A valid \egglog program should not explicitly refer to 
 a specific uninterpreted constant $n$.
We include uninterpreted constants in the syntax nonetheless
 since they are useful when describing the semantics of \egglog programs.

% \rw{Do we need $f(p_1, \ldots, p_k)$ as another case in Atom? I think having $f(p_1, \ldots, p_k) \mapsto 0$ is enough?}

% \newcommand\nt[1]{\ensuremath{\mathit{#1}}}
\begin{figure}
  \centering
  \begin{tabular}{llcl}
     Program        & $P$ & ::= & $R_1,\ldots, R_n$                                                                 \\
     Rule           & $R$ & ::= & $A\ngets A_1,\ldots A_m$.                 \\
     Atom           & $A$ & ::= & $f(p_1,\ldots, p_k)\mapsto o\mid f(p_1,\ldots, p_k)$                 \\
     Pattern        & $p$ & ::= & $f(p_1, \ldots, p_k)\mid o$                 \\
     Term           & $t$ & ::= & $f(t_1, \ldots, t_k)\mid v$                 \\
     Base pattern          & $o$ & ::= & $v\mid x$                                               \\
     Constant       & $v$ & ::= & $c\mid n$ \\
     Interpreted Constant       & $c\in C$                                           \\
     Uninterpreted Constant       & $n\in N$                                           \\
     Variable       & $x, y, \ldots$  \\
  \end{tabular}
  \caption{
    Syntax of core \egglog. 
    % Moreover, in Section \ref{sec:semantics}, 
    % we also use $n\in N$ to denote an uninterpreted constant
    % and $t\in N\cup C$ to denote a constant in general.
    }
  \label{fig:syntax}
\end{figure}

\subsection{Semantics}\label{sec:semantics}

\def\DB{\mathit{DB}}
\def\bmt{\textbf{t}}
\def\bmv{\textbf{v}}
\def\mergef{\mathit{merge}}
Given an infinite set of uninterpreted constants\footnote{
    These uninterpreted constants 
    play a similar role as \eclass ids in \eqsat
    or labelled nulls in the chase from the database literature.
}
 $N = \{n_1,n_2, \ldots\}$ 
 and a complete lattice $L = (C, \sqsubseteq, \sqcup)$ 
 over domain of interpreted constants $C$,
 We define $\bot$ to be the least element of $L$.
A schema $S$ is a collection of function symbols and their function signatures, 
 where the types range over $\{N, C\}$.
Given a schema $S$, an instance of $S$ is defined $I=(\DB, \equiv)$,
 where $\DB$ is a set of function entries $f(v_1,\ldots, v_k)\mapsto v$
 that is consistent with the schema
 and $\equiv$ is an equivalence relation over $N\cup C$ 
 satisfying $\forall c_1, c_2\in C. c_1\equiv c_2\rightarrow c_1=c_2$ 
 (i.e., interpreted constants are only equivalent to themselves).
\revision{For convenience, we also lift set operator (e.g., union, difference)
 to be between an instance and a database, 
 which applies the operator to instance's database.}

Given an arbitrary total order $<$ over $N\cup C$,
we define the canonicalization function $\lambda_{\equiv}(t)=\min t':t'\equiv t$.
For convenience, we lift $\lambda_{\equiv}$ to 
 also work on sets 
 and the whole database instances by pointwise canonicalization.

Before proceeding to define the semantics of an \egglog program,
 we need to first define what it means for a ground atom (an atom without variables)
 to be in the database and what it means to add one to the database.
First, we use the judgement $I\vdash A$ 
 to denote a ground atom is contained in the database $I$.
% \pz{Suggestion: First, we use the judgement $I\vdash A$ 
% to denote a ground atom $A$ is contained in the database $I$. yz: done}
\begin{mathpar}
    \inferrule{
        I\vdash t_i\mapsto v_i \text{ for $i=1\ldots k$}\\\\\and
        I=(\DB{}, \equiv)\and
        f(v_1, \ldots, v_k)\mapsto v \in \DB{}
    }{
        I\vdash f(t_1, \ldots, t_k)\mapsto v
    }
    \and
    \inferrule{
    }{
        I\vdash v\mapsto v
    }
    \and
    \inferrule{
        I\vdash f(t_1, \ldots, t_k)\mapsto v \text{ for some $v$}
    }{
        I\vdash f(t_1, \ldots, t_k)
    }
\end{mathpar}
% We also loose the definition and use $I\vdash_{\lambda} A$ to denote whether an atom is in the database after canonicalization, i.e., $I\vdash_{\lambda} A$ iff $I\vdash \lambda_{\equiv}(A)$.

\def\flatten{\textit{flatten}}
\def\aux{\textit{aux}}

We also define $\flatten{}_I$ in \autoref{fig:flatten}
 to flatten function entries to be inserted 
 into $I$ given a nested ground atom $A$.

\begin{figure}
\begin{align*}
    \flatten{}_I(A) &= s\quad \text{where $(v, s) = \aux{}(A)$}\\
    \aux{}(f(t_1,\ldots, t_k)\mapsto v)
      & = \left(
        v,
        \{f(v_1,\ldots, v_k)\mapsto v\}\cup \bigcup_{i=1,\ldots,k}s_i
      \right)\\
      & \text{where $(v_i, s_i)=\aux{}(t_i)$ for $i=1,\ldots,k$.}\\
    \aux{}(f(t_1,\ldots, t_k))
      & = \left(
        v,
        \{f(v_1,\ldots, v_k)\mapsto v\}\cup \bigcup_{i=1,\ldots,k}s_i
      \right)\\
      & \text{where $(v_i, s_i)=\aux{}(t_i)$ for $i=1,\ldots,k$}\\
      & \text{and $I\vdash f(v_1,\ldots, v_k)\mapsto v$ 
        if such $v$ exists and $v=\textit{default}_f$ otherwise.}\\
    \aux{}(v) &= \left(v, \emptyset\right)
\end{align*}
\caption{
$\flatten{}_I(A)$ 
 flattens function entries to be inserted 
 into $I$ given a nested ground atom $A$.
If the output type of $f$ is $N$,
 then $\textit{default}_f$ is a fresh constant from $N$; 
 otherwise it is $\bot$.
The auxillary function $\aux$ 
 takes a ground atom and
 returns the ``output value'' of the ground atom
 and the set of flattened facts
 it will populate.
}\label{fig:flatten}
\end{figure}
% We say $f(\bmt{})\mapsto t\in I$ 
%  if $\DB{}[f](\lambda_{\equiv}(\bmt{}){})\equiv t$.

Now we can define the semantics of an \egglog program.
It consists of two parts: the immediate consequence operator and the rebuilding operator.
We can define the (inflationary) immediate consequence
 operator $T^\uparrow_P$ as follows.%
 \footnote{
     The definition of immediate consequence operator
      in standard Datalog does not union with $\DB{}$,
     because rule applications in standard Datalog are monotone.
     This is not the case in \egglog in general.
     For example, rule $Q(e)\ngets \textit{lo}(e)\mapsto 5$, where 
      $\textit{lo}$ tracks the lower bound of an expression,
     is not monotone because the value of $\textit{lo}(e)$ can increase over time.
     Although one can adapt the meet semantics of Flix~\citep{flix} for relational joins to enforce monotonicity, 
     we do not do this in \egglog to be compatible with existing \egg applications, 
     which can be non-monotonic.
     Instead, we define \egglog semantics using the inflationary immediate consequence operator, 
     which is used to describe semantics for non-monotonic extension of Datalog such as $\text{Datalog}^\neg$~\citep{neg-by-fixpoint}.
 }
Let $\sigma$ denote a substitution that maps variables to constants,
 and let $A[\sigma]$ denote the ground atom obtained by applying $\sigma$ to atom $A$
 in the standard way.
Given an \egglog program $P$ 
 consisting of a set of rules and $I=(\DB{}, \equiv)$,
 then $T^{\uparrow}_P(I)=\DB\cup T_P(I)$ and $T_P(I)=(\DB{}', \equiv)$, where
\[
    \DB{}'=
    \bigcup_{
        (A\ngets A_1,\ldots, A_m) \in P
    }
    \left\{ 
        \displaystyle
        \textit{flatten}_I(A[\sigma])\mid
        \forall_{i=1,\ldots,m}\ 
        I\vdash A_i[\sigma]
    \right\}
\]

However, functions in $T^\uparrow_P(I)$ may no longer preserve the functional dependencies, 
 as it is possible that the same key $(v_1,\ldots, v_k)$ are mapped to more than one $v$ in some $f$.
We call $T^\uparrow_P(I)$ a \emph{pre-instance}, since it is not a valid instance yet.
To transform a pre-instance into a valid instance, 
 we further define the rebuilding operator \(
    R((\DB{}, \equiv)) = (\DB{}_R, \equiv_R)
\), where:
\begin{align*}
    (\equiv_R) &=
    \text{equivalence closure of }
    \left(
    (\equiv) \cup
    \left\{
        (n_1,n_2)\;
        \bigg|
        \begin{array}{l}
          f(v_1,\ldots, v_k)\mapsto n_1\in \DB{}, \\
          f(v_1,\ldots, v_k)\mapsto n_2\in \DB{}, \\
          n_1, n_2\in N
        \end{array}
    \right\}
    \right)
    \\
    \DB{}_R &=
    \lambda_{\equiv_R}
    \left(\left\{
        f(v_1,\ldots, v_k)
        \mapsto 
        \mergef_{f,\equiv}\left(
            K
        \right)
        \bigg|
        \begin{array}{r}
          K=\left\{v: f(v_1,\ldots, v_k)\mapsto v\in \DB{}\right\} \\
          \text{and $K$ is not empty}
        \end{array}
    \right\}\right)
    \\
    \mergef_{f,\equiv}(K) &=
    \begin{cases}
        \min\left(\lambda_{\equiv}(K)\right) &\text{if the output type of $f$ is $N$;}\\
        \bigsqcup K &\text{if the output type of $f$ is $C$.}
    \end{cases}
\end{align*}
% \[
%     \equiv\cup 
%     \{(n_1,n_2)\mid f(v_1,\ldots, v_k)\mapsto n_1\in \DB{}, 
%     f(v_1,\ldots, v_k)\mapsto n_2\in \DB{}, n_1, n_2\in N\}
% \]
% and
% \[
%     % \DB{}_R[f]=\left\{
%     %     \bmt\mapsto 
%     %     \mergef_{f,\equiv}\left(
%     %         \left\{t: f\left(\bmt\right)\mapsto t\in I\right\}
%     %     \right)\mid
%     %     \text{for all $\bmt$ where $\left\{t: f\left(\bmt\right)\mapsto t\in I\right\}$ is not empty}
%     % \right\}
%     \DB{}_R=\lambda_{\equiv_R}\left(\left\{
%     f(v_1,\ldots, v_k)\mapsto 
%     \mergef_{f,\equiv}\left(
%         K
%     \right)\mid 
%     \text{ $K=\left\{v: f(v_1,\ldots, v_k)\mapsto v\in \DB{}\right\}$ and $K$ is not empty}
%     \right\}\right)
% \]
% where \[
%     \mergef_{f,\equiv}(K)=
%     \min\left(\lambda_{\equiv}(K)\right)
% \]
% if the output type of $f$ is $N$ and \[
%     \mergef_{f,\equiv}(K)=\bigsqcup K
% \] if the output type of $f$ is $C$.
% where
%  \[
%     \mergef_{f,\equiv}(K)=
%     \begin{cases}
%         \min\left(\lambda_{\equiv}(K)\right) &\text{if the output type of $f$ is $N$;}\\
%         \bigsqcup K &\text{if the output type of $f$ is $C$.}
%     \end{cases}
% \]
% \pz{Have we lost the definition of $\lambda_{\equiv_R}$?  Should this be $\mergef_{f,\equiv_R}$? }

\def\lfp{\mathit{lfp}}

Note that the canonicalization in computing $\DB{}_R$ 
 (i.e., $\lambda_{\equiv_R}$) may cause $I$ to be invalid again,
 so successive rounds of rebuilding may be needed. 
Therefore, the complete rebuilding function $R^{\infty}$ 
 is defined as iterative applications of $R$ until fixpoint.
\revision{The rebuilding process always terminates, 
 as it shrinks the size of the database in each iteration.}

We define one iteration of evaluating 
 an \egglog program $F_P$ as $R^{\infty}\circ T^\uparrow_P$, i.e.,
 do rule application once, and apply rebuilding until fixpoint. 
Intuitively, applying $F_P$ to a database makes it ``larger'', 
 in the sense that more facts may be represented.
To capture this monotonicity, we define the expanded database 
 $E_{\equiv}(\DB{})$
 such that 
 $f(v_1, \ldots, v_k)\mapsto n\in E_{\equiv}(\DB{})$ iff
  $f(\lambda_{\equiv}(v_1),\ldots, \lambda_{\equiv}(v_k))\mapsto \lambda_{\equiv}(n)\in \DB{}$
 and
 $f(v_1, \ldots, v_k)\mapsto c\in E_{\equiv}(\DB{})$ iff
  $\exists c'. f(\lambda_{\equiv}(v_1),\ldots, \lambda_{\equiv}(v_k))\mapsto c'\in \DB{}\land c\sqsubseteq c'$.
A database is larger than or equal to another database
 if it knows at least as many facts and equalities,
 so we define
 $(\DB_1, \equiv_1)\sqsubseteq_I (\DB_2, \equiv_2)$ iff
 $E_{\equiv_1}(\DB_1)\subseteq E_{\equiv_2}(\DB_2)$ 
 and
 $\equiv_1\subseteq\equiv_2$.

Although $F_P$ is not a monotone function in general,
the following sequence of iterative applications is always monotonically increasing:
\[
    I_{\bot}\sqsubseteq_I F_P(I_{\bot})\sqsubseteq_I F_P(F_P(I_{\bot}))\sqsubseteq_I  F_P(F_P(F_P(I_{\bot})))\sqsubseteq_I \ldots
\]
for initial database $I_{\bot}=(\emptyset, \textit{Id}_{N\cup C})$
where $\textit{Id}_{N\cup C}$ is the identity relation
over $N\cup C$. This ensures the existence of a fixpoint.

Finally, the result of evaluating an \egglog program $P$ is defined as the inductive fixpoint of $F_P$, 
 i.e., \(
    [\![P]\!]=F_P^{\infty}(I_{\bot})% =\sup\left(\{F_P^n(I_{\bot})\mid n\in\mathbb{N}\}\right)
\).
% If $F_P$ is monotone, then $F_P^{\infty}(I_{\bot})$ is also the least fixpoint of $F_P$.
For many practical applications, 
 $[\![P]\!]$ often has an infinite size, 
 so we only calculate a finite under-approximation of the result, 
 i.e., $(R^{\infty}\circ T^\uparrow_P)^n(I_\bot)$ for some iteration size $n$. 
\revdel{
This gives the \naive algorithm for evaluating an \egglog program given iteration size $n$ shown in Algorithm~\ref{alg:eval}.
}

\def\UF{\textit{UF}}

\revision{
\subsection{Semi-na\"ive Evaluation}\label{sec:seminaive}
}

% We now present a \seminaive algorithm for evaluating an \egglog program.
% Algorithm~\ref{alg:eval} shows the pseudo-code for the algorithm.
% We begin by initializing the database $DB_0$, 
% the differential database $\Delta\DB{}_0$, 
% and the Union-Find $\UF$ to be empty.
% Each iteration of the main loop computes the new differential database $\Delta\DB{}_i$.
% For each rule $A \ngets A_1, \ldots, A_m$ in the program, 
%  $\textsc{evalQuery}$ generates $m$ \emph{delta rules}~\footnote{
% Certain variants of the semi-na\"ive algorithm keeps track of both 
% the ``current'' and ``previous'' versions of the database to 
% further speed up the evaluation.
% For simplicity, we only use the current database.
% }
%  $\{\Delta A \ngets A_1, \ldots, A_{j-1}, \Delta A_j, A_{j+1}, \ldots, A_m 
%  \mid j \in [1\ldots m]\}$, 
%  then evaluates each delta rule to compute the new facts in $\Delta\DB{}_i$.
% Next, we remove the old facts in $DB_{i-1}$ from $\Delta\DB{}_i$, 
%  and combine the old database with the differential into $DB_i$. 
% Finally, we rebuild $DB_i$, $\Delta\DB{}_i$, and $\UF$ to resolve functional dependencies.
% During rebuilding, 
%  any new or updated tuple is also added to $\Delta\DB{}_i$, 
%  and old tuples are removed from both $DB_i$ and $\Delta\DB{}_i$.

\def\SN{\mathit{SN}}
\def\N{\mathit{N}}

Last section gives an algorithm for evaluating \egglog programs, 
 which iteratively apply the immediate consequence operator ($T^\uparrow_P$)
 and the rebuilding operator $R^\infty$.
We call this algorithm the \naive evaluation.
Despite straightforward,
 the \naive evaluation may duplicate works by re-discovering same facts
 over and over again.
The \seminaive algorithm mitigates this problem 
 by incrementalizing the evaluation.
In \seminaive evaluation, each iteration maintains a differential database $\Delta\DB_i$,
 which will only contain tuples that are updated or new in this iteration.
The \seminaive rule application operator $T^\SN_P(I_i, \Delta\DB_i)$ 
 additionally takes a differential database.
For each rule $A\ngets A_1, \ldots, A_m$, $T^\SN_P$ will expand it
 into $m$ delta rules 
 $\{A \ngets A_1, \ldots, A_{j-1}, \Delta A_j, A_{j+1}, \ldots, A_m \mid j \in 1\ldots m\}$
 and apply these rules to the database similar to $T_P$.

\begin{algorithm}
    \caption{The \seminaive evaluation of an \egglog program.}
    \label{alg:eval}
    \begin{algorithmic}
    \Procedure{$F^\SN_P$}{$P, n$}
        \State $I_0\gets I_\bot$;\ \ $\Delta\DB{}_0\gets \emptyset$;
        \For{$i = 1\ldots n$}
            \State $\left(\DB_i{},\equiv_i\right) \gets R^\infty\left(
                I_{i-1}\cup T^\SN_P\left(
                    I_{i-1}, \Delta\DB_{i-1}
                \right)
            \right)$;
            \State $\Delta\DB{}_i \gets \DB{}_{i} - \DB{}_{i-1}$;
            \State $I_i\gets \left(\DB_i{},\equiv_i\right)$;
        \EndFor
        \State \Return $I_n$;
    \EndProcedure
    \end{algorithmic}
\end{algorithm}


% \begin{algorithm}
%     \caption{The \seminaive evaluation of an \egglog program.
%     Removing the underlined parts results in the \naive algorithm.
%      }
%     \label{alg:eval}
%     \begin{algorithmic}
%     \Procedure{EvalEqLog}{$P, n$}
%         \State $\DB{}_0, \Delta\DB{}_0, \UF{} \gets \textsc{newDB}(), \textsc{newDB}(), \textsc{newUnionFind}()$;
%         % \State $\UF{} \gets \textsc{newUnionFind}()$
%         \For{$i = 1\ldots n$}
%             \State $\DB{}_i \gets DB{}_{i-1}$
%             % \State \uwave{$\Delta\DB{}\gets \textsc{newDB}()$};
%             \For{$A\ngets A_1,\ldots, A_m \in P$} \Comment{Search}
%                 % \For{\underline{rule $\in \textsc{deltaRules}$(rule)}}
%                     % \State $A\ngets A_1,\ldots, A_m \gets$ rule;
%                     \For{$\textit{subst}\in \textsc{evalQuery}((A_1,\ldots, A_m), \DB{}_{i-1}, \underline{\Delta\DB{}_{i-1}})$}
%                        \State $\DB{}_i.\textit{add}(
%                            A.\textit{subst}(\textit{subst})
%                        )$;
%                        \Comment{Apply}
%                     \EndFor
%                 % \EndFor
%             \EndFor
%             % \State $\DB{}_i \gets \DB{}_{i-1} \cup \Delta\DB{}_{i}$; \Comment{Apply}
%             % \State $\DB{}_i,\UF{} \gets \textit{rebuild}(\DB{}_i,\UF{})$; \Comment{Rebuild}
%             \Repeat \Comment{Rebuild}
%                 \State $\DB_i{},\UF_i{} \gets \textit{rebuildOnce}(\DB{}_i,\UF_i{})$;
%             \Until{$DB$ does not change};

%             \State $\Delta\DB{}_i \gets \DB{}_{i} - \DB{}_{i-1}$;
%             \If{$\Delta\DB{}_i=\emptyset$}
%                 \State \Return $(\DB{}_i, \UF{})$;
%             \EndIf
%         \EndFor
%         \State \Return $(\DB{}_n, \UF{})$;
%     \EndProcedure
%     \end{algorithmic}
% \end{algorithm}

% The fixpoint semantics in \autoref{sec:semantics} directly prescribes a \emph{na\"ive} evaluation algorithm
%  of \egglog programs.
% It differs from the \seminaive algorithm in three ways:
% first, $\textsc{evalQuery}$ simply runs each rule as-is using only $DB_{i-1}$;
%  second, $\Delta\DB{}_i$ is directly merged with $\DB{}_i$ without removing old facts;
%  finally, rebuilding need only operate on $DB_i$ and not $\Delta\DB{}_i$.
% In other words, removing the underlined parts from Algorithm~\ref{alg:eval} 
% results in the na\"ive algorithm.

We prove the following theorem in the full version \citep{egglog-preprint}.

\begin{theorem}
    The \seminaive evaluation of an \egglog program
    produces the same database as the \naive
    evaluation.
\end{theorem}

% Old proof

% To see the semi-na\"ive algorithm computes the same result as the na\"ive algorithm,
% we need to establish the following two properties: 
% first, the result of $\textsc{evalQuery}$ for the semi-na\"ive algorithm
%  is a subset of its result for the na\"ive algorithm;
% second, any tuple produced by the na\"ive algorithm
%  but not by the semi-na\"ive algorithm
%  must already be in the database.
% The first condition  holds because $\Delta\DB{}_i$ is always a subset of $\DB{}_i$
%  when $\textsc{evalQuery}$ is called, 
%  and since we disallow negated relations, every \emph{query} is monotone in 
%  each relation (even though the \emph{action} may not be monotone).
% For the second condition, we first establish an invariant for $rebuild$: 
% the difference $DB_i - \Delta\DB{}_i$ after rebuilding is a subset 
%  of their difference before rebuilding.
% This is because any tuple added by $rebuild$ is added to both $DB_i$ and $\Delta\DB{}_i$,
%  and any tuple removed from $\Delta\DB{}_i$ is also removed from $DB_i$, 
%  if it is present there. 
% Next, we note that any tuple $t$ produced by na\"ive but not \seminaive 
%  must also be produced by evaluating the query over $\DB_{i-1}-\Delta\DB_{i-1}$.
% Now since $\DB_{i-1}-\Delta\DB_{i-1}$ before rebuilding is itself a subset 
%  of $DB_{i-2}$ (because of the line $DB_i \gets DB_{i-1} \cup \Delta\DB{}_i$), 
%  we have that $t$ must have already been produced by running the query on 
%  $DB_{i-2}$, and therefore is already in the database. 
% Finally, we can show the semi-na\"ive algorithm is equivalent to the na\"ive 
% algorithm by induction on the number of iterations $i$,
% using the two properties above.





% % \pz{Might be interesting to note "finite (under) approximation" to re-emphasize that the sense this is 
% % an approximation is rigorous}
% \def\UF{\textit{UF}}
% \begin{algorithm}
%     \caption{Evaluating an \egglog program}
%     \label{alg:eval}
%     \begin{algorithmic}
%     \Procedure{EvalEqLog}{$P,n$}
%         \State $\DB{}, \UF{} \gets \textsc{newDB}(), \textsc{newUnionFind}()$;
%         % \State $\DB{} \gets \textsc{newDB}()$
%         % \State $\UF{} \gets \textsc{newUnionFind}()$
%         \For{$i=1\ldots n$}
%             \State $\Delta\DB{}\gets \textsc{newDB}()$;
%             \For{$A\ngets A_1,\ldots, A_m \in P$} \Comment{Search}
%                  \For{$\textit{subst}\in \textsc{evalQuery}((A_1,\ldots, A_m), \DB{})$}
%                     \State $\Delta\DB{}.\textit{add}(
%                         A.\textit{subst}(\textit{subst})
%                     )$;
%                  \EndFor
%             \EndFor
%             \State $\DB{}.\textit{insertAll}(\Delta\DB{})$; \Comment{Apply}
%             \If{$\DB{}$ does not change}
%                 \State \textbf{break};
%             \EndIf
%             \Repeat \Comment{Rebuild}
%                 \State $\DB{},\UF{} \gets \textit{rebuildOnce}(\DB{},\UF{})$;
%             \Until{$DB$ does not change};
%         \EndFor
%         \Return $(\DB{}, \UF{})$;
%     \EndProcedure
%     \end{algorithmic}
% \end{algorithm}
